\documentclass[11pt]{article}

\usepackage[margin=1in]{geometry}
\usepackage{amsmath,amssymb}
\usepackage{booktabs}
\usepackage{tabularx}
\usepackage{array}
\usepackage{enumitem}
\usepackage[T1]{fontenc}
\usepackage{hyperref}

\newcommand{\code}[1]{\texttt{#1}}
\newcommand{\bD}{\boldsymbol{\Delta}}
\newcommand{\nb}{\bar n}
\newcolumntype{L}{>{\raggedright\arraybackslash}X}
\emergencystretch=3em

\title{The field pair correlation $g(r)$ in \code{cluster\_diagnostics.py}\\
and \code{cluster\_diagnostics\_3d.py}: definition, normalization and scaling}
\author{FHDeX}
\date{}

\begin{document}
\maketitle

\noindent This note collects the properties of the density-field $g(r)$ written
to \code{rdf\_field.txt} by the 2D and 3D cluster tools: what it measures, why
an ideal gas gives $g=1$, what $g(0)$ means and how it scales with the grid and
the dimension, and how $g(r)$ relates to the probability of finding a particle
in a shell. Everything here holds in both 2D and 3D; only the cell measure and
the shell measure change.

\section{Definition}

The domain is periodic with volume $V$ (an area in 2D), divided into $M$ cells
of volume $\Delta V$, so $V=M\,\Delta V$. From the density $\phi$ (normalized
so that $\sum_i\phi_i\Delta V=1$) and the particle number $N$
(\code{-{}-num-part}) the tools form cell counts
\begin{equation}
  n_i = \phi_i\,N\,\Delta V ,\qquad \sum_i n_i = N ,\qquad \nb = N/M .
\end{equation}
For a lattice displacement $\bD$ between cells (with periodic wrap $i+\bD$),
\begin{equation}
  g(\bD) = \frac{M}{N(N-1)}\sum_i n_i\,n_{i+\bD} \quad(\bD\ne0),\qquad
  g(0) = \frac{M}{N(N-1)}\sum_i n_i\,(n_i-1) .
  \label{eq:def}
\end{equation}
The sums run over all $M$ cells. The code evaluates $\sum_i n_in_{i+\bD}$ for
every $\bD$ at once as the circular autocorrelation of $n$ by FFT, subtracts
$N$ at $\bD=0$, and multiplies by $M/(N(N-1))$. It then bins the
displacements by $|\bD|$ (minimum image) and divides each bin by its exact
number of displacements $c_b$ (the \code{ndisplacements} column), so the value
written for a bin is the average of $g(\bD)$ over the displacements in it:
\begin{equation}
  g_b = \frac1{c_b}\sum_{\bD\in b} g(\bD) .
\end{equation}
With the default bin width $\Delta r=\Delta x$ the first bin, $[0,\Delta x)$,
holds only $\bD=0$, so the first row of \code{rdf\_field.txt} is $g(0)$.
The results are averaged over the frames in \code{-{}-t-range}.

\section{\texorpdfstring{$g(\bD)$ as a pair probability}{g as a pair probability}}

Pick an ordered pair $(a,b)$ of distinct particles at random; there are
$N(N-1)$ such pairs. Let $P(\bD)$ be the probability that $b$'s cell is $a$'s
cell shifted by $\bD$. Counting pairs cell by cell,
\begin{equation}
  P(\bD) = \frac{1}{N(N-1)}\sum_i n_i\,n_{i+\bD}
  \quad\bigl(n_i(n_i-1)\text{ at }\bD=0\bigr),
\end{equation}
because $n_in_{i+\bD}$ is the number of ordered pairs with one particle in cell
$i$ and the other in cell $i+\bD$ (and $n_i(n_i-1)$ the number of ordered pairs
of distinct particles both in cell $i$). Comparing with \eqref{eq:def},
\begin{equation}
  \boxed{\;P(\bD) = \frac{g(\bD)}{M} = g(\bD)\,\frac{\Delta V}{V}\;}
  \label{eq:prob}
\end{equation}
So $g$ is the pair probability per displacement divided by its value $1/M$ for
uncorrelated particles. The probabilities sum to one exactly: summing
$\sum_in_in_{i+\bD}$ over all $\bD$ gives $N^2$, and the self-pair subtraction
removes $N$, so
\begin{equation}
  \frac1M\sum_{\bD} g(\bD) = 1 ,
\end{equation}
the discrete form of $\frac1V\int g\,dV=1$. (The file is restricted to
$r\le L/2$, so it does not contain every displacement.)

\section{\texorpdfstring{Why an ideal gas gives $g=1$}{Why an ideal gas gives g = 1}}

``Uniform'' must mean \emph{statistically} uniform: $N$ particles placed
independently, each uniformly over the domain (an ideal gas, Poisson
statistics). Write $[a\in i]$ for 1 if particle $a$ is in cell $i$ and 0
otherwise, so $n_i=\sum_a[a\in i]$, and each particle is in a given cell with
probability $1/M$, independently of the others.

\paragraph{$\bD\ne0$.}
\begin{equation}
  \sum_i n_in_{i+\bD} = \sum_i\sum_a\sum_b [a\in i]\,[b\in i+\bD] .
\end{equation}
Terms with $a=b$ vanish, since one particle cannot be in two different cells.
For each of the $N(N-1)$ ordered pairs $a\ne b$ the expectation is $1/M^2$, so
\begin{equation}
  \mathbb E\Bigl[\sum_in_in_{i+\bD}\Bigr] = M\cdot\frac{N(N-1)}{M^2}
  = \frac{N(N-1)}{M}
  \quad\Longrightarrow\quad \mathbb E[g(\bD)] = 1 .
\end{equation}

\paragraph{$\bD=0$.} Now $\sum_in_i^2$ also contains the $N$ terms $a=b$
(each particle paired with itself):
\begin{equation}
  \mathbb E\Bigl[\sum_in_i^2\Bigr] = N + \frac{N(N-1)}{M} ,
\end{equation}
and the $-n_i$ in $n_i(n_i-1)$ removes exactly those self pairs, so again
$\mathbb E[g(0)]=1$. Without the subtraction one would get
$g(0)=1+M/(N-1)\approx1+1/\nb$: the extra $1/\nb$ is the Poisson shot noise of
the counts, $\mathrm{Var}(n_i)=\nb$, not a correlation.

A single frame fluctuates around these values; each bin averages many
displacements and the file averages frames, so the measured value is close to
1 (the Poisson check of the tools gives $g=1.00000$). In the SPDE, the
Dean--Kawasaki noise gives an ideal-gas field the same count statistics (mean
and variance $\nb$), so a noisy non-interacting run also gives $g\approx1$ at
all $r$, including $r=0$. The in-code check $S(k)\approx1$ for the ideal gas is
the same statement in $k$ space.

\paragraph{Perfectly equal counts.} If instead $n_i=\nb$ exactly in every cell
(for example a noise-free run from a uniform state),
\begin{equation}
  g(\bD\ne0) = \frac{M\cdot M\nb^2}{N(N-1)} = \frac{N}{N-1}\approx1,\qquad
  g(0) = \frac{M\cdot M\nb(\nb-1)}{N(N-1)} = \frac{N-M}{N-1}\approx 1-\frac1{\nb}.
\end{equation}
With $\nb=5$, $g(0)\approx0.8$. Equal counts have no fluctuation at all
($\mathrm{Var}=0$ instead of $\nb$), so subtracting self pairs removes more than
the fluctuations contributed, and $g(0)$ dips below 1, as for a lattice. A
deterministic run with $g(0)<1$ therefore means a field smoother than an ideal
gas, not repulsion.

\section{\texorpdfstring{What $g(0)$ measures}{What g(0) measures}}

\subsection{Mass-weighted density enhancement}
Rewriting \eqref{eq:def} with weights $w_i=n_i/N$ (the fraction of the
particles in cell $i$) and $\phi_i/\bar\phi=n_i/\nb$,
\begin{equation}
  g(0) = \frac{N}{N-1}\Bigl\langle\frac{\phi}{\bar\phi}\Bigr\rangle_{w}
         - \frac{M}{N-1}
  \;\approx\; \Bigl\langle\frac{\phi}{\bar\phi}\Bigr\rangle_{w} - \frac1{\nb},
  \qquad
  \Bigl\langle\frac{\phi}{\bar\phi}\Bigr\rangle_{w} = \sum_i w_i\,\frac{\phi_i}{\bar\phi}.
\end{equation}
$g(0)$ is the density a typical particle sees in its own cell, relative to the
mean, with the shot noise removed. It is a measure of how concentrated the
mass is \emph{at the scale of one cell}.

\subsection{Limiting cases and the occupied fraction}
\begin{center}
\begin{tabularx}{\textwidth}{@{}L l l@{}}
\toprule
Configuration & occupied fraction of cells & $g(0)$ \\
\midrule
ideal gas (Poisson) & 1 & 1 \\
equal counts $n_i=\nb$ & 1 & $1-1/\nb$ \\
$K$ equal clusters, each uniform over $c$ cells & $Kc/M$ & $\approx M/(Kc)$ \\
$K$ equal clusters, each in a single cell & $K/M$ & $\approx M/K$ \\
all $N$ particles in one cell & $1/M$ & $M$ \\
\bottomrule
\end{tabularx}
\end{center}
For all particles at one point, $n=N$ in one cell and 0 elsewhere, so
$g(0)=M\,N(N-1)/(N(N-1))=M$ and $g(\bD)=0$ for $\bD\ne0$: $2.1\cdot10^6$ at
$128^3$, $1.6\cdot10^4$ at $128^2$. In the continuum this is $g(r)=V\delta(r)$,
which is unbounded; only the grid caps it at $M$. In general,
\begin{equation}
  g(0)\approx\frac{1}{\text{fraction of the domain the mass occupies}} ,
\end{equation}
measured in cells. If a mass fraction $f_c$ is spread uniformly over a
fraction $a$ of the cells and the rest uniformly over the remainder,
\begin{equation}
  g(0)\approx \frac{f_c^2}{a}+\frac{(1-f_c)^2}{1-a} .
\end{equation}
Real clusters are peaked rather than flat, which raises the mass-weighted
density, so this underestimates $g(0)$. For $K$ Gaussian clusters of width
$\sigma\gg\Delta x$, number density $n_c=K/V$ and total mass fraction $f_c$,
\begin{equation}
  g(0)\approx\frac{f_c^2}{n_c\,(4\pi\sigma^2)^{d/2}} ,
\end{equation}
roughly the volume per cluster over the volume of a cluster.

\subsection{Why 3D values are larger: dimension and resolution}
The definition is the same in 2D and 3D; what changes is the measure of a cell
($\Delta x^2$ against $\Delta x^3$) and how much of the domain a cluster fills.
A cluster of radius $a$ in a box of side $L$ fills a fraction $\pi(a/L)^2$ in
2D but $\frac43\pi(a/L)^3$ in 3D: for $a/L=0.02$, $1.3\cdot10^{-3}$ against
$3.4\cdot10^{-5}$, about 37 times smaller. So even many more clusters in 3D can
fill a smaller fraction of the space, and $g(0)$ is correspondingly larger. The
runs in \code{exec/dean\_kow/interaction}, last frames:
\begin{center}
\begin{tabular}{@{}l r r r r r@{}}
\toprule
 & $g(0)$ & $\langle\phi/\bar\phi\rangle_w$ & $\nb$ & 90\% of mass in & estimate \\
\midrule
2D \code{\_new\_repul\_2} & 9.36 & 9.56 & 5 & 12.9\% of cells & 7.9 \\
3D \code{\_3d\_noise} & 62.0 & 62.2 & 9.5 & 2.6\% of cells & 31 \\
\bottomrule
\end{tabular}
\end{center}
$g(0)$ equals $\langle\phi/\bar\phi\rangle_w-1/\nb$ in both. The estimate uses
flat clusters: 2D, 53 clusters of about 20 cells ($a=6.5\%$) holding
$f_c=0.71$ of the mass; 3D, 476 clusters of about 120 cells ($a=2.7\%$) holding
$f_c=0.92$. The measured values are higher because the clusters are peaked.
(The first row of \code{rdf\_field.txt} for the 2D run reads 7.05 because it
averages all frames, including the uniform $t=0$.)

\paragraph{Resolution dependence.} $g(0)$ is a cell average, not the continuum
limit $g(r\to0)$. It grows under grid refinement whenever the clusters are
only a few cells across (the 3D clusters, about 3 cells in radius, are in that
regime). To compare runs, use a resolution-independent quantity such as the
cluster size \code{Rg} in \code{clusters.txt} or the volume fraction inside
clusters; $g(0)$ says how compact the clusters are, not how they are arranged
(for that, see \code{order\_diagnostics.tex}).

\subsection{\texorpdfstring{Relation to $S(k)$}{Relation to S(k)}}
With $\hat n(\mathbf k)=\sum_in_ie^{-i\mathbf k\cdot\mathbf x_i}$ and
$S(\mathbf k)=|\hat n|^2/N$, Parseval's theorem gives
$\sum_in_i^2=\frac1M\sum_{\mathbf k}|\hat n(\mathbf k)|^2
=\frac1M\bigl(N^2+N\sum_{\mathbf k\ne0}S(\mathbf k)\bigr)$, hence exactly
\begin{equation}
  g(0) = 1+\frac{1}{N-1}\sum_{\mathbf k\ne0}\bigl(S(\mathbf k)-1\bigr)
  \;\approx\; 1+\frac{1}{\nb}\,\overline{\bigl(S-1\bigr)} ,
\end{equation}
where the sum runs over all $M-1$ nonzero wave vectors of the FFT grid,
including the corner modes beyond the Nyquist circle, and the overline is the
mean over them. Summing the complete shells in \code{sf.txt} leaves those
corner modes out and does not reproduce $g(0)$ exactly.

\section{Shell probabilities in 2D and 3D}

The continuum statements are
\begin{equation}
  \text{2D: } P(\text{shell})\approx g(r)\,\frac{2\pi r\,\Delta r}{V},\qquad
  \text{3D: } P(\text{shell})\approx g(r)\,\frac{4\pi r^2\,\Delta r}{V},
\end{equation}
where $P(\text{shell})$ is the probability that one particular other particle
lies in the shell $[r,r+\Delta r)$ around a given particle; the expected number
of neighbours in the shell is $(N-1)P=\rho\,g(r)\times(\text{shell measure})$
with $\rho=(N-1)/V$. The field $g$ satisfies this exactly on the lattice.
Summing \eqref{eq:prob} over the displacements of bin $b$,
\begin{equation}
  P(\text{bin }b) = \sum_{\bD\in b}g(\bD)\,\frac{\Delta V}{V}
  = g_b\,\frac{c_b\,\Delta V}{V},
\end{equation}
exactly, where $c_b\,\Delta V$ is the lattice version of the shell measure:
the number of cell displacements in the shell times the cell measure. For
$r\gg\Delta x$ it tends to the continuum value,
\begin{center}
\begin{tabular}{@{}l l l@{}}
\toprule
 & continuum shell & lattice $c_b\,\Delta V$ \\
\midrule
2D & $2\pi r\,\Delta r$ & $\#\{(i,j): r\le|(i\Delta x,j\Delta y)|<r+\Delta r\}\times\Delta x\Delta y$ \\
3D & $4\pi r^2\Delta r$ & $\#\{(i,j,k)\text{ in the shell}\}\times\Delta x\Delta y\Delta z$ \\
\bottomrule
\end{tabular}
\end{center}
which gives the 2D and 3D statements above. At small $r$ the lattice is coarse
and the counts differ: the bin $[\Delta x,2\Delta x)$ holds 8 displacements in
2D (the 4 at distance $\Delta x$ and the 4 at $\sqrt2\Delta x$) against
$\pi(2^2-1^2)\approx9.4$ cells for the continuum annulus, and 26 in 3D (6, 12
and 8 at $1$, $\sqrt2$, $\sqrt3$ times $\Delta x$) against
$\frac43\pi(2^3-1)\approx29.3$. Because the code divides by the exact count
$c_b$ rather than the continuum formula, each $g_b$ is an unbiased average and
the identity holds exactly at every $r$; only the shortcut ``$2\pi r$'' or
``$4\pi r^2$'' is approximate.

\paragraph{The cluster-centre $g(r)$} (\code{rdf\_clusters.txt}) works the
other way round and uses the continuum shell measure directly. For $n_c$
centres and pairs counted in $[r_{\mathrm{lo}},r_{\mathrm{hi}})$,
\begin{equation}
  g = \frac{\text{pair count}}{\frac12n_c(n_c-1)\,\text{shell}/V},\qquad
  \text{shell}=\begin{cases}
    \pi\,(r_{\mathrm{hi}}^2-r_{\mathrm{lo}}^2) & \text{2D},\\[2pt]
    \frac43\pi\,(r_{\mathrm{hi}}^3-r_{\mathrm{lo}}^3) & \text{3D},
  \end{cases}
\end{equation}
accumulated over frames. The denominator is the expected pair count for
uniformly random centres, so $P(\text{a given other centre in the shell})=
g\,\text{shell}/V$ holds by construction.

\paragraph{Cell-centre smearing.} The field $g$ locates particles only to
their cells, so each displacement is between cell centres: $g_b$ is the true
pair correlation smeared over about $\pm\Delta x$ in each direction. This
matters for $r$ of a few $\Delta x$, the same regime where $c_b\Delta V$ and
the continuum shell measure differ.

\end{document}
